%------------------------------------------------------------------------------%
\begin{question}
  Calcule o vetor gradiente da função \(f(x, y) = \sqrt{2x+3y}\)
  no ponto \((-1, 2)\)
\end{question}

%------------------------------------------------------------------------------%
\begin{resolution}{Calculando a Derivada Parcial em $x$}
  \[
    \D{f}{x}
    \pause
    \;=\; \D{}{x}\sqrt{2x+3y}
    \pause
    \;=\; \D{}{x}\left(2x+3y\right)^{\sfrac{1}{2}}
    \pause
    \;=\; \frac{1}{2}\left(2x+3y\right)^{\sfrac{-1}{2}}\D{}{x}\left(2x+3y\right)
  \]

  \[
    \hphantom{\D{f}{x}}
    \pause
    \;=\; \frac{1}{2}\,\frac{1}{\sqrt{2x+3y}}\, 2\D{x}{x}
    \pause
    \;=\; \frac{1}{\sqrt{2x+3y}}
  \]

  \pause
  \bigskip
  \[
    \D{f}{x}(-1, 2)
    \pause
    \;=\; \frac{1}{\sqrt{2(-1)+3(2)}}
    \pause
    \;=\; \frac{1}{\sqrt{-2+6}}
    \pause
    \;=\; \frac{1}{\sqrt{4}}
    \pause
    \;=\; \frac{1}{2}
  \]
\end{resolution}

%------------------------------------------------------------------------------%
\begin{resolution}{Calculando a Derivada Parcial em $y$}
  \[
    \D{f}{y}
    \pause
    \;=\; \D{}{y}\sqrt{2x+3y}
    \pause
    \;=\; \D{}{y}\left(2x+3y\right)^{\sfrac{1}{2}}
    \pause
    \;=\; \frac{1}{2}\left(2x+3y\right)^{\sfrac{-1}{2}}\D{}{y}\left(2x+3y\right)
  \]

  \[
    \hphantom{\D{f}{y}}
    \pause
    \;=\; \frac{1}{2}\,\frac{1}{\sqrt{2x+3y}}\, 3\D{y}{y}
    \pause
    \;=\; \frac{3}{2\sqrt{2x+3y}}
  \]

  \pause
  \bigskip
  \[
    \D{f}{y}(-1, 2)
    \pause
    \;=\; \frac{3}{2\sqrt{2(-1)+3(2)}}
    \pause
    \;=\; \frac{3}{2\sqrt{-2+6}}
    \pause
    \;=\; \frac{3}{2\sqrt{4}}
    \pause
    \;=\; \frac{3}{4}
  \]
\end{resolution}

%------------------------------------------------------------------------------%
\begin{resolution}{O Vetor Gradiente}
  \[
    \nabla f(x, y)
    \pause
    \;=\; \Vetor{\D{f}{x}}{\D{f}{y}}
    \pause
    \;=\; \Vetor{\frac{1}{\sqrt{2x+3y}}}{\frac{3}{2\sqrt{2x+3y}}}
    \pause
    \;=\; \frac{1}{\sqrt{2x+3y}} \vetor{1}{\sfrac{3}{2}}
  \]
\end{resolution}

%------------------------------------------------------------------------------%
\begin{resolution}{O Vetor Gradiente no Ponto \((-1, 2)\)}
\begin{align*}
    \onslide<+->{\nabla f(-1, 2)}
    \onslide<+->{& = \left[\frac{1}{\sqrt{2x+3y}} \vetor{1}{\sfrac{3}{2}}\right]\evalat{(-1, 2)}{} \\}
    \onslide<+->{& = \frac{1}{\sqrt{2(-1)+3\times 2}} \vetor{1}{\sfrac{3}{2}} \\}
    \onslide<+->{& = \frac{1}{\sqrt{4}} \vetor{1}{\sfrac{3}{2}} \\}
    \onslide<+->{& = \vetor{\sfrac{1}{2}}{\sfrac{3}{4}} }
\end{align*}
\end{resolution}

%------------------------------------------------------------------------------%
